\begin{textsample}{POS Dim 3 – rj – Score 26.00 – rj/rj\_inf\_natureza\_2.txt}  \label{ex:f3_pos_005}
Mas como é possível questionar e intervir nos espaços educacionais em busca da \textbf{instituição} de novas formas de viver e de pensar a vida na Terra?
Buscando respostas, foram apontadas proposições organizadas em torno de três objetivos : ( 1 ) o primeiro desafia uma cultura antropocêntrica que entende os seres humanos como autônomos, superiores e, portanto, donos do destino das demais espécies e de tudo o que existe ; ( 2 ) o segundo busca reinventar os caminhos de conhecer, numa perspectiva que considere as múltiplas dimensões do humano ; e ( 3 ) o terceiro \textbf{convida} a dizer não ao consumismo e ao desperdício de recursos naturais ( TIRIBA, 2007 ).
A - Religar as \textbf{crianças} com a natureza : desemparedar Em 1988, quando foi aprovada a atual Constituição Brasileira, a Educação \textbf{Infantil} passou a ser um direito das \textbf{crianças}.
Mas, se elas chegam às \textbf{Instituições} de Educação \textbf{Infantil} aos quatro \textbf{meses} e saem aos cinco anos ; se, até os dois anos, frequentam raramente o \textbf{pátio}, e, a partir dessa idade, adquirem o direito de permanecer por apenas uma ou duas horas ao ar livre, \textbf{brincando} sobre cimento, brita ou grama sintética ; se as janelas da sala onde permanecem o restante do tempo não permitem a visão do mundo exterior ; se assim os \textbf{dias} se sucedem, essas \textbf{crianças} não conhecem a liberdade...
Assim, \textbf{cuidar} das \textbf{crianças} significa mantê -las em contato com o universo natural de que são parte.
Se o nosso compromisso é com a sua integridade e com a preservação da vida no planeta, Sol, ar puro, água, terra, barro, \textbf{areia} são elementos/condições que devem estar presentes no \textbf{dia} a \textbf{dia} de \textbf{creches} e \textbf{pré-escolas}.
Os \textbf{bebês} vão gostar muito de estar ao ar livre, sobre colchonetes, desfrutando do espaço aberto, \textbf{atentos} ao que está ao redor.
As \textbf{crianças} de dois e três anos poderão passear no entorno da escola, acompanhadas pelas turmas maiores, que adoram \textbf{cuidar}, \textbf{brincar} com eles, \textbf{conversar}.
Nesta mesma linha de raciocínio, podemos pensar que as \textbf{brincadeiras} nos espaços externos podem constituir fonte de sentimentos de solidariedade e companheirismo.
Um \textbf{pátio} que é de todos, onde cada um pode escolher com quem e com que deseja \textbf{brincar}, não favorece atitudes individualistas e competitivas ; ao contrário, constitui espaço de convivência amistosa, prazerosa.
B- Reinventar os caminhos de conhecer Conhecer é sentir, pois todo sistema racional tem um fundamento emocional ( MATURANA, 2002 ).
Antes de lidar com conceitos abstratos, as \textbf{crianças} deveriam aprender a apreciar e a amar um lugar ( ORR, 1995 ).
Pode ser um \textbf{pequeno} vale, as margens de um riacho, com seus pássaros ; um manguezal, uma montanha, uma praia, onde céu, nuvens, ventos, animais compõem um cenário de \textbf{brincadeiras} e descobertas, constituindo -se, a partir daí, como objeto de investigação pedagógica.
É a possibilidade de estar neste lugar que pode promover o encontro com aquilo que verdadeiramente importa a cada \textbf{criança} ou ao grupo e, portanto, será capaz de mantê -las interessadas ( FREINET, 1979 ).
Um rio, por exemplo : de onde vem e para onde corre, que seres o habitam ou utilizam suas águas, quem vive às suas margens?
Partindo de uma relação com uma realidade ecológica concreta, as respostas a essas perguntas são encontradas na Biologia, na Geografia, na História, na Sociologia, etc., exigindo, portanto, uma pesquisa pedagógica que não pode deter -se nessa ou naquela área de estudos, mas atravessa e interconecta infinitos campos do conhecimento, é transdisciplinar ( ALVES, GARCIA, 2001 ; GALLO, 2001 ).
Portanto, não se trata de aprender o que é uma árvore decompondo -a em suas partes.
Mas de senti -la e compreendê -la em interação com a vegetação que está ao redor, com os animais que se alimentam de seus frutos, com as nuvens que trazem chuva, com a sensação \textbf{agradável} gerada pela sombra em que \textbf{brincamos}.
Experiências de plantio de hortaliças, flores e ervas e temperos possibilitam às \textbf{crianças} essa \textbf{percepção} ecológica da realidade, em que as interações entre seres, coisas e fenômenos tendem sempre para um todo coerente e complexo ( MATURANA, VARELA, 2002 ).
Valorizando atividades de plantar, \textbf{colher} e comer alimentos sem agrotóxicos, estaremos abrindo espaços para o exercício da ética do \textbf{cuidado} em relação ao próprio corpo, à Terra, ao entorno, ao planeta.
Mas essas experiências não podem ser eventuais, devem estar no coração do projeto pedagógico, como \textbf{rotina}, de tal forma que as \textbf{crianças} tenham acesso direto e \textbf{frequente}, reguem, participem da limpeza da horta, da colheita, integrando -se, vivenciando e conhecendo, na prática, os processos de \textbf{nascimento} e \textbf{crescimento} dos frutos da terra.
Isso nada tem a ver com as experiências em que as \textbf{crianças} “ plantam ” feijão sobre o algodão molhado no copinho e depois que ele brota jogam tudo no lixo.
Se abandonarmos o minhocário depois que as \textbf{crianças} entendem a importância da minhoca no trato agrícola ; se deixamos sem água as mudas recém-brotadas, se mantivermos em cativeiro os animais tão comuns nos \textbf{pátios} das escolas, como porquinho-da-índia e jabuti, ensinaremos a \textbf{meninos} e \textbf{meninas} uma visão utilitarista da natureza, atitudes de desrespeito a seres vivos.

% matched lemmas: agradável, areia, atento, bebê, brincadeira, brincar, colher, conversar, convidar, creche, crescimento, criança, cuidado, cuidar, dia, frequente, infantil, instituição, menino, mês, nascimento, pequeno, percepção, pré-escola, pátio, rotina
\end{textsample}
